\scanpage{175}
\mspage{490}
% position: upper long paragraph
jours. Puis, entraîné par le petit diable (ou ange, je ne sais...) de la sys-\linebreak tématique en moi, j'ai fini par ressortir mes \inserted{vieilles réflexions d'il y a} cinq ans sur ce thème. Pendant une semaine ou deux je m'étais amusé alors à ``ramasser'' une centaine ou deux de ces couples bien suggestifs, lesquels s'étaient alors assemblés par affinités en une vingtaine de groupes. Comme cette réflexion s'est faite en marge du fameux ``ouvrage poétique'' que j'étais en train d'écrire, je n'ai pu m'empêcher de ranger ces groupes tant bien que mal à la queue-leu-leu, par affinités et filiations de sens d'un groupe \inserted{au sui-}\linebreak vant. Hier soir, reprenant la réflexion avec du recul, et sans carcan poétique autour du cou, c'est dix-huit groupes que j'ai trouvés(au lieu de vingt), par un groupage peut-être un peu plus rigoureux. Je soupçonne d'ailleurs qu'il doit y avoir bien d'autres groupes encore, peut-être même un nombre illimité, corres-\linebreak pondants à des modes d'appréhension de la réalité auxquels je n'ai pas songé au cours du travail (ni, peut-être, jamais encore).

% position: middle long paragraph
Quant aux dix-huit groupes que j'ai bel et bien relevés, je me suis effor-\linebreak cé de les assembler en un diagramme (ou ``graphe'') suivant les principaux liens d'affinités qui les relient les uns aux autres. Certains de ces liens d'ailleurs ne se sont imposés à mon attention qu'au cours du tracé d'ébauches successives du diagramme. Le travail ici était vraiment très proche du travail mathémati-\linebreak que bien familier, quand on s'efforce de saisir graphiquement, de façon aussi frappante que possible, un ensemble plus ou moins complexe de relations (don-\linebreak nées par exemple par des ``applications'', figurées par des flèches) entre un certain nombre d'``ensembles'' ou de ``catégories'', figurant comme ``sommets'' du ``diagramme'' qu'on s'efforce de construire. Là aussi, des exigences de nature essentiellement esthétiques, de symétrie et de transparence structurale notam-\linebreak ment, conduisent fréquemment à introduire (et au besoin donc, à découvrir voire même à inventer) des ``flèches'' ou liens auxquels on n'avait pas songé au départ, et parfois même des nouveaux ``sommets''. Toujours est-il qu'après cinq ou six ébauches successives, j'ai fini par aboutir à un diagramme, vaguement en forme d'arbre de Noël, qui m'a satisfait provisoirement - d'autant plus qu'il commen-\inserted{\unsure{:}}\linebreak çait vraiment à se faire prohibitivement tard!

% position: lower short paragraph
Je me suis couché content, je sentais que je n'avais pas perdu mon temps, même si nes notes n'avaient pas avancé d'un poil (*). Mais je m'étais remis en
% position: footnote rule
\par\noindent\rule{0.25\linewidth}{0.4pt}
% position: footnote
(*) En compensation, je pourrais déposer un brevet sur l'invention d'une nou-\linebreak velle forme \inserted{poétique, savoir le poème dit ``non linéaire'', ou ``diagramma-\linebreak tique''.}
% position: foot pencil
\margin{174}
